\documentclass[12pt, a4paper]{report}

% Packages
\usepackage[utf8]{inputenc}
\usepackage[T1]{fontenc}
\usepackage[english]{babel}
\usepackage{amsmath, amssymb, amsthm, amsfonts}
\usepackage{geometry}
\usepackage{hyperref}
\usepackage{bookmark}
\usepackage{enumitem}
\usepackage{fancyhdr}
\usepackage{titlesec}
\usepackage{xcolor}
\usepackage{tikz}
\usetikzlibrary{trees}

% Geometry setup
\geometry{
    top=2.5cm,
    bottom=2.5cm,
    left=2.5cm,
    right=2.5cm
}
\setlength{\headheight}{30pt}

% Hyperref setup
\hypersetup{
    colorlinks=true,
    linkcolor=blue,
    filecolor=magenta,      
    urlcolor=cyan,
    pdftitle={Stochastic Finance and Risk Modelling - Chapter 2},
    pdfauthor={Laurence Carassus and Gaoyue Guo},
}

% Theorem environments
\newtheorem{theorem}{Theorem}[chapter]
\newtheorem{proposition}[theorem]{Proposition}
\newtheorem{lemma}[theorem]{Lemma}
\newtheorem{corollary}[theorem]{Corollary}

\theoremstyle{definition}
\newtheorem{definition}{Definition}[chapter]
\newtheorem{example}{Example}[chapter]
\newtheorem{remark}[theorem]{Remark}

% Layout/Style
\pagestyle{fancy}
\fancyhf{}
\rhead{\leftmark}
\lhead{Chapter 2: Stochastic Processes}
\cfoot{\thepage}

\title{\textbf{Stochastic Finance and Risk Modelling} \\ \large Chapter 2: From Probability to Stochastic Process (Continuation)}
\author{Based on slides by Ioane Muni-Toke \\ Adapted by Laurence Carassus and Gaoyue Guo \\ \textit{Expanded Study Resource}}
\date{\today}

\begin{document}

\maketitle

\tableofcontents

\chapter*{Introduction}
\addcontentsline{toc}{chapter}{Introduction}

This document serves as a comprehensive study guide for the second part of the "From probability to stochastic process" lecture. It expands upon the key concepts of stochastic processes, filtrations, stopping times, and martingales in discrete time.

The goal is to provide a rigorous yet intuitive understanding of how information evolves over time in a probabilistic setting, which is the cornerstone of financial modeling. We will explore the mathematical structures that allow us to model random systems that evolve, culminating in the powerful theory of martingales.

\chapter{Stochastic Processes}

\section{Definition and Basic Properties}

A stochastic process is a mathematical object defined to describe the evolution of a system subject to randomness over time. Formally, it is a collection of random variables defined on a common probability space.

\begin{definition}[Stochastic Process]
Let $(\Omega, \mathcal{F}, \mathbb{P})$ be a probability space. A \textbf{stochastic process} is a family of random variables $X = (X_t)_{t \in I}$ defined on $(\Omega, \mathcal{F})$ and taking values in a measurable space $(E, \mathcal{E})$, often called the \textit{state space}.
\end{definition}

The set $I$ is called the \textbf{index set} or \textbf{parameter set}, representing time.
\begin{itemize}
    \item \textbf{Discrete-time}: If $I$ is a countable set, such as $\mathbb{N} = \{0, 1, 2, \dots\}$ or a finite set $I = \{0, 1, \dots, T\}$. In this course, we primarily focus on the discrete-time setting with a finite horizon $I = \{0, 1, \dots, T\}$.
    \item \textbf{Continuous-time}: If $I$ is an interval of $\mathbb{R}$, such as $[0, T]$ or $[0, \infty)$. This is the subject of advanced stochastic calculus (e.g., Brownian motion).
\end{itemize}

The state space $E$ is typically $\mathbb{R}^d$ equipped with the Borel $\sigma$-algebra $\mathcal{B}(\mathbb{R}^d)$.

\section{The Canonical Construction}

One standard way to construct a stochastic process is to define it on a \textit{product space}. This is often referred to as the "canonical" construction, where the sample space $\Omega$ itself represents the set of all possible paths of the process.

Consider a finite horizon $T \in \mathbb{N}$ and a state space $E$.
\begin{itemize}
    \item Let $\Omega = E^T = \{ \omega = (\omega_1, \dots, \omega_T) : \omega_i \in E \}$ be the set of all possible sequences of length $T$.
    \item We can define a $\sigma$-algebra $\mathcal{F}$ on $\Omega$, typically the product $\sigma$-algebra.
    \item We define a probability measure $\mathbb{P}$ on $(\Omega, \mathcal{F})$.
\end{itemize}

In this setting, the stochastic process $X = (X_t)_{1 \leq t \leq T}$ can be defined as the \textbf{coordinate mapping}:
\[
X_t(\omega) = \omega_t, \quad \text{for } \omega = (\omega_1, \dots, \omega_T).
\]
Here, the randomness $\omega$ \textit{is} the entire trajectory of the process, and $X_t$ simply "reads" the value of the trajectory at time $t$.

\subsection{Example: Coin Toss (Random Walk)}

Consider the "Head or Tail" game where we toss a balanced coin $T$ times.
\begin{itemize}
    \item \textbf{State Space}: $E = \{H, T\}$ (Head, Tail). Alternatively, we can map this to numerical values, e.g., $\{1, 0\}$ or $\{1, -1\}$.
    \item \textbf{Sample Space}: $\Omega = E^T = \{ \omega = (\omega_1, \dots, \omega_T) : \omega_i \in \{H, T\} \}$. The size of the sample space is $|\Omega| = 2^T$.
    \item \textbf{Probability Measure}: Since the coin is balanced, every outcome sequence is equally likely.
    \[
    \mathbb{P}(\{\omega\}) = \frac{1}{|\Omega|} = \frac{1}{2^T}, \quad \forall \omega \in \Omega.
    \]
    The $\sigma$-algebra is the power set $\mathcal{F} = 2^\Omega$.
\end{itemize}

We can define random variables $X_t$ representing the outcome of the $t$-th toss:
\[
X_t(\omega) = \begin{cases} 1 & \text{if } \omega_t = H \\ 0 & \text{if } \omega_t = T \end{cases}
\]
The collection $(X_t)_{1 \leq t \leq T}$ is a stochastic process taking values in $\{0, 1\}$.

We can also define the \textbf{accumulated sum} process $S = (S_t)_{1 \leq t \leq T}$:
\[
S_t = \sum_{i=1}^t X_i
\]
The process $(S_t)$ takes values in the set $\{0, 1, \dots, T\}$ and represents the total number of "Heads" obtained up to time $t$.

\section{Trajectories (Sample Paths) and Notation}

The most complete and rigorous way to define a stochastic process is as a function of two variables: time and randomness.
\[
X : I \times \Omega \to E
\]
Let us break down this structural notation:
\begin{itemize}
    \item $X$: The name of the stochastic process.
    \item $I$: The \textbf{index set} (time domain), e.g., $\{0, 1, \dots, T\}$.
    \item $\Omega$: The \textbf{sample space} (the set of all possible random outcomes $\omega$).
    \item $\times$: The \textbf{Cartesian product}, meaning $X$ takes a pair of inputs $(t, \omega)$.
    \item $\to E$: The \textbf{state space}. The $\to$ arrow defines the structural mapping from the domain sets to the codomain set $E$.
\end{itemize}

When we write $X(t, \omega) = x$, we mean: "At time $t$, given that the random scenario $\omega$ occurred, the value of the process is $x \in E$."

This two-variable perspective allows us to "fix" one variable and study the process from two complementary views. To understand these views, it is crucial to distinguish between two types of arrows used in mathematical mapping:
\begin{itemize}
    \item The straight arrow \textbf{$\to$} (e.g., $A \to B$) defines the \textit{sets} involved: it maps a domain to a codomain.
    \item The maps-to arrow \textbf{$\mapsto$} (e.g., $x \mapsto f(x)$) defines the \textit{operation rule}: it shows exactly how a specific element from the domain is transformed into an element in the codomain.
\end{itemize}

\subsection{View 1: A collection of random variables (Fixing time)}
If we fix a specific time $t \in I$ and let $\omega$ vary, we obtain a new function:
\[
X(t, \cdot) : \Omega \to E, \quad \omega \mapsto X_t(\omega)
\]
Here, $X(t, \cdot)$ is a \textbf{Random Variable}. The dot "$\cdot$" indicates that the randomness variable is left free. We usually drop the dot and simply write \textbf{$X_t$}. It answers the question: "What will the value be at time $t$?" The answer is a random distribution over $E$.

\subsection{View 2: A random function / Trajectory (Fixing randomness)}
If the random outcome has already been realized (we fix a specific scenario $\omega \in \Omega$) and we let time run, we obtain:
\[
X(\cdot, \omega) : I \to E, \quad t \mapsto X_t(\omega)
\]
For a fixed outcome $\omega$, the deterministic function $t \mapsto X_t(\omega)$ is called a \textbf{trajectory}, \textbf{sample path}, or \textbf{realization} of the process.

\begin{definition}[Trajectory]
For each $\omega \in \Omega$, the function $X(\cdot, \omega)$ describes one specific history of the system's evolution over time.
\end{definition}

If we denote $\Sigma = E$ (the state space) and $\Sigma^I$ as the set of all functions from $I$ to $\Sigma$, then the entire process $X$ can be viewed as a single random variable taking values in the function space $\Sigma^I$:
\[
X: (\Omega, \mathcal{F}) \to (\Sigma^I, \mathcal{G}^{\otimes I})
\]
where $\mathcal{G}^{\otimes I}$ is the approproiate product $\sigma$-algebra on the function space.

\chapter{Filtrations and Information}

In financial modeling, and stochastic processes in general, the concept of \textit{information} is crucial. We need a mathematical way to represent "what is known" at a given time $t$. This is the role of a filtration.

\section{Filtrations}

\begin{definition}[Filtration]
A \textbf{filtration} on a probability space $(\Omega, \mathcal{F}, \mathbb{P})$ is an increasing family of sub-$\sigma$-algebras $(\mathcal{F}_t)_{t \in I}$ of $\mathcal{F}$. That is:
\[
\mathcal{F}_s \subseteq \mathcal{F}_t \subseteq \mathcal{F}, \quad \forall s \le t, \quad s, t \in I.
\]
A probability space equipped with a filtration, $(\Omega, \mathcal{F}, (\mathcal{F}_t)_{t \in I}, \mathbb{P})$, is called a \textbf{filtered probability space}.
\end{definition}

\subsection{Interpretation}
Each $\sigma$-algebra $\mathcal{F}_t$ represents the collection of events whose occurrence (or non-occurrence) is known by time $t$. The condition $\mathcal{F}_s \subseteq \mathcal{F}_t$ for $s \le t$ reflects the fact that information \textit{accumulates} over time: we do not forget what we knew in the past.
\begin{itemize}
    \item If $A \in \mathcal{F}_t$, then at time $t$, an observer knows whether the event $A$ has happened or not.
    \item $\mathcal{F}_0$ represents the initial information (often trivial, containing only $\emptyset$ and $\Omega$, if the process hasn't started or no prior information is available).
    \item $\mathcal{F}_T$ (or $\mathcal{F}_{\infty}$) represents the total information available at the end of the timeline.
\end{itemize}

\subsection{Example: The Coin Toss Filtration}
Recall our coin toss experiment with $\Omega = \{H, T\}^T$.
\begin{itemize}
    \item Let $\mathcal{F}_0 = \{ \emptyset, \Omega \}$. At the start, we know nothing about the outcome.
    \item Let $\mathcal{F}_t = \sigma(X_1, \dots, X_t)$. This is the $\sigma$-algebra generated by the first $t$ tosses.
\end{itemize}
Information at time $t$ consists of knowing the precise sequence of Heads and Tails up to step $t$, but nothing about steps $t+1, \dots, T$.
For $T=3$ and $t=1$:
\[
\mathcal{F}_1 = \sigma(X_1) = \{ \emptyset, \{H \times \{H,T\}^2\}, \{T \times \{H,T\}^2\}, \Omega \}
\]
Knowing which set in $\mathcal{F}_1$ contains $\omega$ is equivalent to knowing the value of $X_1(\omega)$.

\section{Adapted and Predictable Processes}

The relationship between a stochastic process and a filtration is fundamental. It tells us whether the value of the process is "knowable" at a certain time.

\subsection{Adapted Processes}
\begin{definition}[Adapted Process]
A stochastic process $X = (X_t)_{t \in I}$ is \textbf{adapted} with respect to the filtration $\mathbb{F} = (\mathcal{F}_t)_{t \in I}$ if closely linked to the filtration. specifically:
\[
X_t \text{ is } \mathcal{F}_t\text{-measurable for all } t \in I.
\]
\end{definition}

\textbf{Interpretation}: An adapted process extends the idea of "non-anticipative" behavior. The value $X_t$ is revealed at time $t$. We do not need to look into the future (time $u > t$) to know $X_t$. Historic stock prices are an adapted process: today's price is known today.

\begin{proposition}
If $X$ is adapted, then $X_s$ is $\mathcal{F}_t$-measurable for all $s \le t$.
\end{proposition}
\begin{proof}
Since $X$ is adapted, $X_s$ is $\mathcal{F}_s$-measurable. By definition of filtration, if $s \le t$, then $\mathcal{F}_s \subseteq \mathcal{F}_t$. Therefore, any function measurable with respect to $\mathcal{F}_s$ is automatically measurable with respect to the larger $\sigma$-algebra $\mathcal{F}_t$.
\end{proof}

\subsection{Predictable Processes}
\begin{definition}[Predictable Process]
A stochastic process $X = (X_t)_{t \in I}$ is \textbf{predictable} with respect to the filtration $\mathbb{F} = (\mathcal{F}_t)_{t \in I}$ if:
\[
X_{t} \text{ is } \mathcal{F}_{t-1}\text{-measurable for all } t \in I, \, t \ge 1.
\]
(Typically defining $X_0$ to be $\mathcal{F}_0$-measurable or constant).
\end{definition}

\textbf{Interpretation}: A predictable process is determined by the information available strictly \textit{before} time $t$. Its value at time $t$ is already "locked in" or "decied upon" at time $t-1$.
\begin{itemize}
    \item \textbf{Example}: Trading strategies. If you decide at time $t-1$ to hold $\phi_t$ shares of a stock during the interval $(t-1, t]$, this decision must differ on information available at $t-1$. Thus, the number of shares held is a predictable process.
    \item Note: In discrete time, predictability is a disjoint concept from purely adapted. Every predictable process is adapted (since $\mathcal{F}_{t-1} \subseteq \mathcal{F}_t$), but not every adapted process is predictable (e.g., the stock price itself is adapted but generally not predictable).
\end{itemize}

\section{Generated Filtration}

Every stochastic process generates its own "natural" flow of information.

\begin{definition}[Generated Filtration]
The \textbf{natural filtration} (or generated filtration) of a process $X$, denoted by $\mathbb{F}^X = (\mathcal{F}^X_t)_{t \in I}$, is defined as:
\[
\mathcal{F}^X_t := \sigma(X_s : s \le t) = \sigma(X_0, X_1, \dots, X_t).
\]
\end{definition}

\begin{proposition}
The generated filtration $\mathbb{F}^X$ is the \textbf{smallest} filtration with respect to which $X$ is adapted.
\end{proposition}
\begin{proof}
1. \textbf{Adaptedness}: By definition, $X_t$ is measurable with respect to $\sigma(X_0, \dots, X_t) = \mathcal{F}^X_t$. Thus $X$ is $\mathbb{F}^X$-adapted.
2. \textbf{Minimality}: Let $\mathbb{G} = (\mathcal{G}_t)$ be any other filtration such that $X$ is $\mathbb{G}$-adapted. Then for any $t$, $X_0, \dots, X_t$ are all $\mathcal{G}_t$-measurable (since $X_s \in \mathcal{G}_s \subseteq \mathcal{G}_t$). Since $\mathcal{F}^X_t$ is the \textit{smallest} $\sigma$-algebra making $X_0, \dots, X_t$ measurable, it follows that $\mathcal{F}^X_t \subseteq \mathcal{G}_t$.
\end{proof}

Currently, we often work with the natural filtration unless specified otherwise (e.g., if there are multiple sources of randomness or "insider information").

\chapter{Stopping Times}

In many financial applications (e.g., American options, optimal stopping problems), actions depend on random events. A fundamental concept is the \textbf{stopping time}, which models a decision to stop or take action based \textit{only} on currently available information.

\section{Definition and Intuition}

A random variable $\tau : \Omega \to I \cup \{+\infty\}$ is a time index that is random. However, for it to be useful in decision making, observing the process up to time $t$ must be sufficient to determine if $\tau$ has occurred by time $t$.

\begin{definition}[Stopping Time]
A random variable $\tau$ taking values in $I \cup \{+\infty\}$ is a \textbf{stopping time} with respect to the filtration $\mathbb{F} = (\mathcal{F}_t)_{t \in I}$ if for every $t \in I$:
\[
\{\tau \le t\} \in \mathcal{F}_t.
\]
\end{definition}

\textbf{Intuition}: The event $\{\tau \le t\}$ means "the stopping time has occurred by time $t$". The definition requires this event to be observable given the information $\mathcal{F}_t$.
If $\tau$ represents the time to sell a stock, the decision to sell \textit{at or before time $t$} must rely only on stock prices up to time $t$, not on future prices.

\begin{proposition}
If $I = \mathbb{N}$ (discrete time), then $\tau$ is a stopping time if and only if:
\[
\{\tau = t\} \in \mathcal{F}_t, \quad \forall t \in I.
\]
\end{proposition}
\begin{proof}
$(\Rightarrow)$ If $\tau$ is a stopping time, then $\{\tau \le t\} \in \mathcal{F}_t$ and $\{\tau \le t-1\} \in \mathcal{F}_{t-1} \subseteq \mathcal{F}_t$.
Since $\{\tau = t\} = \{\tau \le t\} \setminus \{\tau \le t-1\}$, it is the difference of two sets in $\mathcal{F}_t$, hence in $\mathcal{F}_t$.
$(\Leftarrow)$ If $\{\tau = k\} \in \mathcal{F}_k$ for all $k$, then
\[
\{\tau \le t\} = \bigcup_{k=0}^t \{\tau = k\}.
\]
Since $\{\tau = k\} \in \mathcal{F}_k \subseteq \mathcal{F}_t$ for $k \le t$, the finite union is in $\mathcal{F}_t$.
\end{proof}

\section{Examples of Stopping Times}

\subsection{First Hitting Time}
Let $A$ be a Borel set in the state space $E$ (e.g., $A = [K, \infty)$ for a price threshold).
Basically, the \textbf{first hitting time} of $A$ is defined as:
\[
H_A(\omega) = \inf \{ t \in I : X_t(\omega) \in A \}.
\]
(Convention: $\inf \emptyset = +\infty$).

\begin{proposition}
If $X$ is an adapted process, then $H_A$ is a stopping time.
\end{proposition}
\begin{proof}
We check the condition $\{\tau \le t\} \in \mathcal{F}_t$:
\begin{align*}
\{H_A \le t\} &= \{ \exists k \in \{0, \dots, t\} \text{ s.t. } X_k \in A \} \\
&= \bigcup_{k=0}^t \{ X_k \in A \}.
\end{align*}
Since $X$ is adapted, $\{ X_k \in A \} \in \mathcal{F}_k \subseteq \mathcal{F}_t$ for all $k \le t$.
Since $\mathcal{F}_t$ is a $\sigma$-algebra, it is closed under finite unions. Thus, $\{H_A \le t\} \in \mathcal{F}_t$.
\end{proof}

\subsection{Counter-Example: Last Hitting Time}
Define $L_A = \sup \{ t \in I : X_t \in A \}$. This is generally \textbf{not} a stopping time.
To know if $L_A \le t$, we need to know that the process will \textit{never} visit $A$ again after time $t$. This requires future information, which is not available in $\mathcal{F}_t$.

\section{Stopped Processes}

Given a process $X$ and a stopping time $\tau$, we often want to consider the process "frozen" at time $\tau$. This is crucial for analyzing strategies that exit a market or stop a game.

\begin{definition}[Stopped Process]
The \textbf{stopped process} $X^\tau = (X^\tau_t)_{t \in I}$ is defined by:
\[
X^\tau_t(\omega) = X_{t \wedge \tau(\omega)}(\omega) = \begin{cases}
X_t(\omega) & \text{if } t < \tau(\omega) \\
X_{\tau(\omega)}(\omega) & \text{if } t \ge \tau(\omega)
\end{cases}
\]
where $a \wedge b = \min(a, b)$.
\end{definition}

\begin{proposition}
If $X$ is adapted and $\tau$ is a stopping time, then the stopped process $X^\tau$ is adapted.
\end{proposition}
\begin{proof}
We need to show that $X^\tau_t$ is $\mathcal{F}_t$-measurable for each $t$.
We can write:
\[
X^\tau_t = X_t \cdot \mathbb{I}_{\{t < \tau\}} + \sum_{k=0}^t X_k \cdot \mathbb{I}_{\{\tau = k\}}.
\]
Let's analyze each term:
1. $\{t < \tau\} = \{\tau \le t\}^c \in \mathcal{F}_t$ (since $\tau$ is a stopping time). $X_t$ is $\mathcal{F}_t$-measurable. Thus the product is $\mathcal{F}_t$-measurable.
2. For each $k \le t$, regarding the term in the sum:
   $\{\tau = k\} \in \mathcal{F}_k \subseteq \mathcal{F}_t$.
   $X_k$ is $\mathcal{F}_k$-measurable $\implies \mathcal{F}_t$-measurable.
   Thus $X_k \cdot \mathbb{I}_{\{\tau = k\}}$ is $\mathcal{F}_t$-measurable.
Sum of measurable functions is measurable.
Therefore, $X^\tau_t$ is $\mathcal{F}_t$-measurable.
\end{proof}

\chapter{Martingales}

The theory of martingales is central to modern probability and finance. A martingale represents the mathematical formalization of a "fair game", where knowledge of the past does not help predict the future mean variations.

\section{Definitions}

Let $(\Omega, \mathcal{F}, (\mathcal{F}_t)_{t \in I}, \mathbb{P})$ be a filtered probability space. Let $X = (X_t)_{t \in I}$ be a stochastic process.
We assume $X$ is \textbf{adapted} and \textbf{integrable}, meaning $X_t \in L^1(\Omega)$ (i.e., $\mathbb{E}[|X_t|] < \infty$) for all $t \in I$.

\begin{definition}[Martingale, Sub/Super-martingale]
The process $X$ is called a:
\begin{enumerate}
    \item \textbf{Martingale} if for all $s \le t$:
    \[
    \mathbb{E}[X_t \mid \mathcal{F}_s] = X_s \quad \text{a.s.}
    \]
    \item \textbf{Sub-martingale} if for all $s \le t$:
    \[
    \mathbb{E}[X_t \mid \mathcal{F}_s] \ge X_s \quad \text{a.s.}
    \]
    \item \textbf{Super-martingale} if for all $s \le t$:
    \[
    \mathbb{E}[X_t \mid \mathcal{F}_s] \le X_s \quad \text{a.s.}
    \]
\end{enumerate}
\end{definition}

\textbf{Intuition}:
\begin{itemize}
    \item \textbf{Martingale}: A "fair game". Given the history up to time $s$, the expected value of your wealth at any future time $t$ is exactly your current wealth $X_s$. There is no drift.
    \item \textbf{Sub-martingale}: A "favorable game" (or representing a convex function of a martingale, like wealth in an economy with growth). The expected future value is greater than or equal to the current value. It tends to grow on average.
    \item \textbf{Super-martingale}: An "unfavorable game" (e.g., gambling at a casino). The expected future value decreases or stays the same.
\end{itemize}

\begin{proposition}
If $X$ is a martingale, then expectations are constant:
\[
\mathbb{E}[X_t] = \mathbb{E}[X_0] = \text{const}, \quad \forall t \in I.
\]
\end{proposition}
\begin{proof}
Take expectations on both sides of the martingale property with $s=0$:
\[
\mathbb{E}[ \mathbb{E}[X_t \mid \mathcal{F}_0] ] = \mathbb{E}[X_0].
\]
By the tower property of conditional expectation, the LHS is $\mathbb{E}[X_t]$.
\end{proof}

\section{Doob's Inequalities}

Martingales allow us to bound the probability of large deviations of the entire path of the process, not just at a fixed time.

\subsection{Doob's Maximal Inequality}

\begin{theorem}[Doob's Sub-martingale Inequality]
Let $X$ be a non-negative sub-martingale. Then for any $\lambda > 0$ and any
$n \in \mathbb{N}$:
\[
\mathbb{P}\left( \max_{0 \le k \le n} X_k \ge \lambda \right) \le \frac{\mathbb{E}[X_n \mathbb{I}_{\{ \max_{k \le n} X_k \ge \lambda \}}]}{\lambda} \le \frac{\mathbb{E}[X_n]}{\lambda}.
\]
\end{theorem}

\begin{proof}
Let $A = \{ \max_{0 \le k \le n} X_k \ge \lambda \}$. We want to bound $\mathbb{P}(A)$.
Define the stopping time $\tau = \inf \{ k : X_k \ge \lambda \} \wedge n$.
Note that on the set $A$, we necessarily have $X_{\tau} \ge \lambda$ (if the maximum is attained before $n$, $X_{\tau} \ge \lambda$; if it is never attained, we don't care about the set $A^c$).
However, let's be more precise.
Let $E_k = \{ \max_{0 \le j < k} X_j < \lambda \text{ and } X_k \ge \lambda \}$. These sets are disjoint and their union is $A = \bigcup_{k=0}^n E_k$.
On the set $E_k$, $X_k \ge \lambda$.
Since $X$ is a sub-martingale, $\mathbb{E}[X_n \mid \mathcal{F}_k] \ge X_k$.
Integrating over $E_k \in \mathcal{F}_k$:
\[
\mathbb{E}[X_n \mathbb{I}_{E_k}] = \mathbb{E}[ \mathbb{E}[X_n \mid \mathcal{F}_k] \mathbb{I}_{E_k} ] \ge \mathbb{E}[ X_k \mathbb{I}_{E_k} ] \ge \lambda \mathbb{P}(E_k).
\]
Summing over $k=0$ to $n$:
\[
\mathbb{E}[X_n \mathbb{I}_A] = \sum_{k=0}^n \mathbb{E}[X_n \mathbb{I}_{E_k}] \ge \lambda \sum_{k=0}^n \mathbb{P}(E_k) = \lambda \mathbb{P}(A).
\]
Thus $\mathbb{P}(A) \le \frac{1}{\lambda} \mathbb{E}[X_n \mathbb{I}_A] \le \frac{1}{\lambda} \mathbb{E}[X_n]$ (since $X_n \ge 0$).
\end{proof}

\subsection{\texorpdfstring{$L^p$}{Lp} Inequality}

A useful consequence for $p > 1$:

\begin{theorem}[Doob's $L^p$ Inequality]
If $X$ is a martingale (or non-negative sub-martingale) and $X_n \in L^p$ with $p > 1$,
then:
\[
\left\| \max_{0 \le k \le n} |X_k| \right\|_p \le \frac{p}{p-1} \|X_n\|_p.
\]
i.e.,
\[
\mathbb{E}\left[ \left( \max_{0 \le k \le n} |X_k| \right)^p \right] \le \left( \frac{p}{p-1} \right)^p \mathbb{E}[|X_n|^p].
\]
\end{theorem}
This theorem essentially says that the maximum of a martingale is controlled by its terminal value in $L^p$ norm.

\section{Doob's Martingale Convergence Theorem}

One of the most powerful results in probability theory concerns the limit behavior of martingales.

\begin{theorem}[Doob's Convergence Theorem]
Let $(X_t)_{t \ge 0}$ be a super-martingale bounded in $L^1$, i.e.,
\[
\sup_{t \ge 0} \mathbb{E}[|X_t|] < \infty.
\]
Then there exists a random variable $X_\infty \in L^1(\Omega)$ such that:
\[
X_t \to X_\infty \quad \text{almost surely as } t \to \infty.
\]
\end{theorem}

\begin{remark}
\begin{itemize}
    \item The condition $\sup \mathbb{E}[|X_t|] < \infty$ is automatically satisfied if $X$ is a non-negative super-martingale (since $\mathbb{E}[|X_t|] = \mathbb{E}[X_t] \le \mathbb{E}[X_0]$).
    \item Convergence is \textbf{almost sure}, not necessarily in $L^1$. For $L^1$ convergence, we need the collection $\{X_t\}$ to be \textit{uniformly integrable}.
\end{itemize}
\end{remark}

\subsection{Interpretation using Upcrossing Inequality}
The proof relies on counting the number of "upcrossings" of an interval $[a, b]$. A convergent sequence cannot oscillate infinitely many times between distinct values $a$ and $b$. Doob showed that the expected number of time a sub-martingale crosses from below $a$ to above $b$ is finite.
\[
\mathbb{E}[U_n[a,b]] \le \frac{\mathbb{E}[(X_n - a)^+]}{b-a}.
\]
Since the RHS is bounded (by the $L^1$ bound), the number of upcrossings is finite almost surely for all rational intervals, implying convergence.


\end{document}
